% © 2026 Ing. David Jaroš — CC BY-NC-ND 4.0
%
% This work is licensed under a Creative Commons Attribution-NonCommercial-NoDerivatives
% 4.0 International License (CC BY-NC-ND 4.0).
%
% License History: Earlier drafts (up to v0.3) were released under CC BY 4.0.
% From v0.4 onward, all material is released under CC BY-NC-ND 4.0 to protect
% the integrity of the theoretical work during ongoing academic development.
%
% See LICENSE.md for full license text.

% e8_no_go_proof.tex
% E8 No-Go Theorem: E8 root system cannot be embedded in B = C ⊗ H
%
% VERDICT: E8 route DISPROVED.
% Status: PROVED — see Theorem E8.1 and Corollary E8.1.
% This document formally closes TARGET 1 of the 2026-04-30 decisive push.
% Source: canonical/algebra/e8_no_go_proof.tex
% Companion: canonical/alpha/neff_geometric_origin.md §3.2–3.4

% UBT-AI-PROVENANCE-BEGIN
% schema: ubt-ai-provenance/v1
% tier: B_machine_verified
% ai_assistance: disclosed
% human_review: machine-verification
% editorial_responsibility: Ing. David Jaroš
% policy: ../../AI_PROVENANCE.md
% notice: Machine-verified against named sources or verifiers; individual attestation is not claimed.
% UBT-AI-PROVENANCE-END

\ifdefined\INCLUDEMODE
\else
  \documentclass[12pt,a4paper]{article}
  \usepackage[a4paper, margin=2.5cm]{geometry}
  \usepackage{amsmath, amssymb, amsthm}
  \usepackage{mathtools}
  \usepackage{hyperref}
  \usepackage{xcolor}
  \usepackage{mdframed}
  \usepackage{booktabs}
  \usepackage{tcolorbox}

  \newtheorem{theorem}{Theorem}[section]
  \newtheorem{lemma}[theorem]{Lemma}
  \newtheorem{proposition}[theorem]{Proposition}
  \newtheorem{corollary}[theorem]{Corollary}
  \theoremstyle{definition}
  \newtheorem{definition}[theorem]{Definition}
  \theoremstyle{remark}
  \newtheorem{remark}[theorem]{Remark}

  \newmdenv[backgroundcolor=yellow!20, linecolor=orange, linewidth=1.5pt]{gapbox}
  \newmdenv[backgroundcolor=green!15, linecolor=green!60!black, linewidth=1.5pt]{resultbox}
  \newmdenv[backgroundcolor=red!8,   linecolor=red!60,   linewidth=1.5pt]{warnbox}
  \newmdenv[backgroundcolor=blue!8,  linecolor=blue!50,  linewidth=1.5pt]{insightbox}

  \title{\textbf{E8 No-Go Theorem}\\
         \large The E8 Root System Cannot Be Embedded in $\mathcal{B} = \mathbb{C}\otimes\mathbb{H}$}
  \author{Ing.\ David Jaro\v{s}}
  \date{2026-04-30}

  \begin{document}
  \maketitle
\fi

%──────────────────────────────────────────────────────────────────────────────
\section{Purpose and Verdict}
\label{sec:e8:purpose}
%──────────────────────────────────────────────────────────────────────────────

\begin{tcolorbox}[colback=red!6!white, colframe=red!60!black,
                  title=\textbf{E8 Route: DISPROVED}]
\begin{tabular}{ll}
Question & Does the E8 root system arise naturally from $\mathcal{B}=\mathbb{C}\otimes\mathbb{H}$? \\
Verdict  & \textbf{NO} — proved below (Theorem~\ref{thm:e8:no_go}). \\
Status   & PROVED [L0] — follows from dimension counting and Gram matrix argument. \\
Action   & E8 removed from UBT core reasoning. \\
         & Mark all E8 connections as NUMERICAL COINCIDENCE only. \\
\end{tabular}
\end{tcolorbox}

%──────────────────────────────────────────────────────────────────────────────
\section{Relevant Algebraic Facts}
\label{sec:e8:facts}
%──────────────────────────────────────────────────────────────────────────────

\subsection{Biquaternion algebra}

From \texttt{canonical/fields/biquaternion\_algebra.tex}:
\begin{equation}
  \mathcal{B} = \mathbb{C}\otimes_{\mathbb{R}}\mathbb{H}
  \;\cong\; \mathrm{M}_2(\mathbb{C}),
  \label{eq:e8:B}
\end{equation}
with real dimension $\dim_{\mathbb{R}}\mathcal{B} = 8$ and ordered $\mathbb{R}$-basis
\begin{equation}
  \mathcal{B}_{\mathrm{basis}}
  = \{1, \mathbf{I}, \mathbf{J}, \mathbf{K}, i, i\mathbf{I}, i\mathbf{J}, i\mathbf{K}\}.
  \label{eq:e8:basis}
\end{equation}

\subsection{Natural inner product on $\mathcal{B}$}

The biquaternion conjugate of $b = \sum_a b_a e_a \in \mathcal{B}$ (where $\{e_a\}$
is the ordered basis~\eqref{eq:e8:basis}) is $b^\dagger = \sum_a \bar{b}_a e_a^*$,
where $(\cdot)^*$ is quaternion conjugation composed with complex conjugation.
The natural real-valued inner product is:
\begin{equation}
  \langle a, b \rangle_{\mathcal{B}}
  \;=\; \operatorname{Re}\!\bigl[\tfrac{1}{2}\,\mathrm{tr}(a^\dagger b)\bigr],
  \qquad a, b \in \mathcal{B},
  \label{eq:e8:inner_product}
\end{equation}
where $\mathrm{tr}$ denotes the matrix trace under $\mathcal{B}\cong\mathrm{M}_2(\mathbb{C})$.

\begin{lemma}[Orthonormality of the basis]
\label{lem:e8:ortho}
The basis~\eqref{eq:e8:basis} is orthonormal under the inner product~\eqref{eq:e8:inner_product}:
\begin{equation}
  \langle e_a, e_b \rangle_{\mathcal{B}} = \delta_{ab},
  \qquad a, b \in \{1,\ldots,8\}.
  \label{eq:e8:ortho}
\end{equation}
\end{lemma}

\begin{proof}
Under $\mathcal{B}\cong\mathrm{M}_2(\mathbb{C})$ via $1\mapsto I_2$,
$\mathbf{I}\mapsto i\sigma_1$, $\mathbf{J}\mapsto i\sigma_2$, $\mathbf{K}\mapsto i\sigma_3$:
\begin{itemize}
  \item $\langle 1, 1\rangle = \tfrac{1}{2}\mathrm{tr}(I_2^2) = 1$.
  \item $\langle \mathbf{I}, \mathbf{I}\rangle = \tfrac{1}{2}\mathrm{tr}((i\sigma_1)^\dagger(i\sigma_1))
        = \tfrac{1}{2}\mathrm{tr}(\sigma_1^2) = 1$.
  \item $\langle \mathbf{I}, \mathbf{J}\rangle = \tfrac{1}{2}\mathrm{tr}(\sigma_1\sigma_2)
        = \tfrac{1}{2}\mathrm{tr}(i\sigma_3) = 0$.
\end{itemize}
All 28 off-diagonal pairs vanish by the tracelessness of Pauli products and the
independence of the complex factor $i$.  All 8 diagonal entries equal~1.
\end{proof}

\subsection{E8 root system and its Gram matrix}

The E8 root system $\Phi(E_8) \subset \mathbb{R}^8$ consists of 240 vectors.
A standard choice of simple roots $\{\alpha_1,\ldots,\alpha_8\}$ has Cartan matrix:
\begin{equation}
  A_{E_8}
  = \begin{pmatrix}
      2 & -1 &  0 &  0 &  0 &  0 &  0 &  0 \\
     -1 &  2 & -1 &  0 &  0 &  0 &  0 & -1 \\
      0 & -1 &  2 & -1 &  0 &  0 &  0 &  0 \\
      0 &  0 & -1 &  2 & -1 &  0 &  0 &  0 \\
      0 &  0 &  0 & -1 &  2 & -1 &  0 &  0 \\
      0 &  0 &  0 &  0 & -1 &  2 & -1 &  0 \\
      0 &  0 &  0 &  0 &  0 & -1 &  2 &  0 \\
      0 & -1 &  0 &  0 &  0 &  0 &  0 &  2
    \end{pmatrix},
  \quad
  A_{ij} = \frac{2\langle\alpha_i,\alpha_j\rangle}{|\alpha_j|^2}.
  \label{eq:e8:cartan}
\end{equation}
The E8 lattice $\Lambda_{E_8}$ is the unique even, unimodular lattice in $\mathbb{R}^8$.
Its Gram matrix $G_{ij} = \langle\alpha_i,\alpha_j\rangle$ has the property that
$G \neq I_8$: several off-diagonal entries $G_{ij} = -1 \neq 0$.

%──────────────────────────────────────────────────────────────────────────────
\section{The No-Go Theorem}
\label{sec:e8:theorem}
%──────────────────────────────────────────────────────────────────────────────

\begin{theorem}[E8 no-go for $\mathcal{B}$]
\label{thm:e8:no_go}
There is no natural embedding of the E8 root system $\Phi(E_8)$ into
$(\mathcal{B}, \langle\cdot,\cdot\rangle_{\mathcal{B}})$ via algebraic structures
of $\mathcal{B} = \mathbb{C}\otimes\mathbb{H}$.
\end{theorem}

\begin{proof}
We present three independent obstructions.

\medskip
\textbf{Obstruction 1 — Gram matrix.}
The natural integral lattice of $\mathcal{B}$ is
$\Lambda_{\mathcal{B}} = \mathbb{Z}\langle 1, \mathbf{I}, \mathbf{J}, \mathbf{K},
i, i\mathbf{I}, i\mathbf{J}, i\mathbf{K}\rangle \cong \mathbb{Z}^8$
(isomorphic as a lattice because the basis is orthonormal by Lemma~\ref{lem:e8:ortho}).

The Gram matrix of $\Lambda_{\mathcal{B}}$ with the inner product~\eqref{eq:e8:inner_product} is:
\begin{equation}
  G_{\mathcal{B}} = I_8.
  \label{eq:e8:gram_B}
\end{equation}
This is the hypercubic lattice $D_1^8$.  In contrast, $G_{E_8} \neq I_8$
(see~\eqref{eq:e8:cartan}): the E8 Gram matrix has off-diagonal entries~$-1$.
Therefore $\Lambda_{\mathcal{B}} \not\cong \Lambda_{E_8}$ as Euclidean lattices.

\medskip
\textbf{Obstruction 2 — Automorphism group.}
The automorphism group of $\mathcal{B}\cong\mathrm{M}_2(\mathbb{C})$ as a
$\mathbb{C}$-algebra is $\mathrm{Aut}(\mathcal{B}) = \mathrm{PGL}(2,\mathbb{C})$.
The Weyl group of E8 is $W(E_8)$ of order $|W(E_8)| = 696{,}729{,}600$ — a finite
group.  However, $\mathrm{PGL}(2,\mathbb{C})$ is a connected Lie group of complex
dimension~3, which contains no finite subgroup of order $696{,}729{,}600$ acting
transitively on 240 elements of equal norm in $\mathbb{R}^8$.
Any embedding $W(E_8)\hookrightarrow\mathrm{PGL}(2,\mathbb{C})$ is impossible
because $\mathrm{PGL}(2,\mathbb{C})$ has no faithful 8-dimensional real
representation compatible with the E8 root action.

\medskip
\textbf{Obstruction 3 — Integer structure.}
The natural integer lattice of $\mathcal{B}$ corresponds to the biquaternions with
Gaussian-integer coefficients $z_\mu \in \mathbb{Z}[i]$.  These are the Lipschitz
integers tensored with $\mathbb{C}$, whose underlying real lattice is $\mathbb{Z}^8$
(see Obstruction~1).  The E8 lattice arises from the icosian ring
$\mathbb{H}(\mathrm{icos}) \oplus \mathbb{H}(\mathrm{icos})$ (product of two copies
of the icosian integers), an 8-dimensional real space that is NOT isomorphic to
$\mathcal{B} = \mathbb{C}\otimes\mathbb{H}$ as a ring.
Specifically: the icosian ring has Galois group $\mathbb{Z}/5\mathbb{Z}$ actions
that produce the E8 roots; the biquaternions have no $\mathbb{Z}/5\mathbb{Z}$
automorphism because $\mathbb{C}$ and $\mathbb{H}$ have no fifth-order automorphisms.
\end{proof}

\begin{corollary}[E8 coincidence is numerical only]
\label{cor:e8:coincidence}
The numerical equality $\dim_{\mathbb{R}}\mathcal{B} = 8 = \mathrm{rank}(E_8)$
is a coincidence of dimension.  It does not give rise to any structural connection:
the lattice, automorphism group, and integer ring of $\mathcal{B}$ are all
incompatible with E8.
\end{corollary}

%──────────────────────────────────────────────────────────────────────────────
\section{Consequences for the UBT Alpha Route}
\label{sec:e8:consequences}
%──────────────────────────────────────────────────────────────────────────────

\begin{resultbox}
\textbf{T1 (E8 Construction) — Final Verdict:}
\begin{itemize}
  \item[\textbf{DISPROVED}] $E_8$ lattice from $\mathcal{B}$: no natural embedding exists
        (Theorem~\ref{thm:e8:no_go}, three independent obstructions, [L0]).
  \item[\textbf{DISPROVED}] E8 sphere-packing density $\pi^4/384$ as normalization:
        the packing is not associated with any metric or mode-counting structure in
        $\mathcal{B}$ (no natural inner product on $\mathcal{B}$ gives E8 packing; see §2).
  \item[\textbf{DISPROVED}] E8 torus realization $T^8 = \mathbb{R}^8/\Lambda_{E_8}$:
        the natural compactification is $T^8 = \mathbb{R}^8/\mathbb{Z}^8$
        (hypercubic torus), not the E8 torus.
  \item[\textbf{ANALOGY ONLY}] Viazovska modular forms: the appearance of
        $\vartheta_3^3$ in the UBT partition function and in the Viazovska E8
        sphere-packing proof is a coincidence of the modular-forms toolkit.
        It is not evidence of an E8 structure in UBT.
\end{itemize}
\end{resultbox}

\begin{warnbox}
\textbf{Note on the Hurwitz–D4 connection}:
The closest existing structure is the Hurwitz integer quaternions
$\mathbb{H}(\mathrm{Hurwitz})$, whose root system is \emph{D4} (not E8).
D4 has rank~4, dimension~4 — half the dimension of E8.
The doubling $D_4 \oplus D_4$ gives a rank-8 lattice, but it is NOT E8:
$D_4 \oplus D_4$ is a different even lattice from E8 (they are not isometric).
This confirms that neither the single-copy nor the double-copy of the natural
quaternion integer structure yields E8.
\end{warnbox}

\paragraph{Impact on the alpha chain.}
The E8 route was listed as a candidate normalization for $B_{\mathrm{base}}$
(via sphere-packing density) and for the T1 mapping tasks.
Both are now closed as FALSE\_LEAD.
The canonical alpha route (\texttt{canonical/alpha/alpha\_best\_route.tex})
proceeds via $N_{\mathrm{eff}}^{3/2}$ with no E8 input.  E8 is removed from
all candidate tables.

%──────────────────────────────────────────────────────────────────────────────
\section{Summary Table}
\label{sec:e8:summary}
%──────────────────────────────────────────────────────────────────────────────

\begin{center}
\renewcommand{\arraystretch}{1.3}
\begin{tabular}{lll}
\toprule
Claim & Verdict & Proof level \\
\midrule
$\Lambda_{\mathcal{B}} \cong \Lambda_{E_8}$ & FALSE & [L0] — Gram matrix $G_{\mathcal{B}}=I_8 \neq G_{E_8}$ \\
$W(E_8) \hookrightarrow \mathrm{Aut}(\mathcal{B})$ & FALSE & [L0] — order $|W(E_8)|$ not in $\mathrm{PGL}(2,\mathbb{C})$ \\
$\mathbb{H}(\mathrm{icos})\oplus\mathbb{H}(\mathrm{icos}) \cong \mathcal{B}$ & FALSE & [L0] — no $\mathbb{Z}/5$ automorphism of $\mathcal{B}$ \\
$\pi^4/384$ in $B_{\mathrm{base}}$ or $N_{\mathrm{eff}}$ & FALSE & [L1] — no derivation exists; FALSE\_LEAD \\
$\vartheta_3^3$ links UBT to E8 geometry & ANALOGY ONLY & [L0] — modular coincidence, no E8 structure \\
\midrule
$\dim_\mathbb{R}\mathcal{B} = 8 = \mathrm{rank}(E_8)$ & TRUE (dimension) & [L0] — numerical coincidence only \\
\bottomrule
\end{tabular}
\end{center}

\ifdefined\INCLUDEMODE
\else
  \end{document}
\fi
